\documentclass[preprint,12pt]{elsarticle}
\usepackage{lmodern}
\usepackage{amssymb,amsmath,mathtools,booktabs,amsthm}
\usepackage[hidelinks]{hyperref}
\newtheorem{theorem}{Theorem}[section]
\newtheorem{proposition}[theorem]{Proposition}
\newtheorem{lemma}[theorem]{Lemma}
\newtheorem{corollary}[theorem]{Corollary}
\theoremstyle{definition}
\newtheorem{definition}[theorem]{Definition}
\newtheorem{remark}[theorem]{Remark}
\newcommand{\F}{\mathbb{F}}
\newcommand{\Z}{\mathbb{Z}}
\newcommand{\Tr}{\operatorname{Tr}}
\newcommand{\GR}{\operatorname{GR}}
\newcommand{\red}{\operatorname{red}}
\newcommand{\pending}[1]{\textbf{[#1]}}

\journal{Finite Fields and Their Applications}
\begin{document}
\begin{frontmatter}
\title{Explicit reduced APN permutations: ramification and optimal Galois-ring lifts}
\author[a,b]{Hiroki Fukui\corref{cor1}}
\ead{fukui@somec.org}
\cortext[cor1]{Corresponding author. ORCID: 0009-0008-7122-522X.}
\affiliation[a]{organization={Research Institute of Criminal Psychiatry / Sex Offender Medical Center},
  addressline={4F Akasaka K-Tower, 1-2-7 Moto-Akasaka}, city={Minato-ku, Tokyo}, country={Japan}}
\affiliation[b]{organization={Department of Neuropsychiatry, Kyoto University}, city={Kyoto}, country={Japan}}
\begin{abstract}
Let $m\ge5$ be odd and $q=2^m$. For every odd $r$ with $3\le r<m$ and $\gcd(r,m)=1$ we give the reduced polynomial $P_{m,r}=\red\bigl((T_r^3)^{2^{m-2}}\bigr)$, $T_r=\sum_{i<r}X^{2^i}$, explicitly: it has coefficients in $\F_2$ and exactly $3r-2$ terms, it is an APN permutation of $\F_q$, and its formal derivative $1+X^{2^{r-2}}+X^{2^{m-2}}$ has no zero on $\F_q$ and is two-to-one. We determine the ramification of $P_{m,r}$ as a map of the projective line for all admissible $(m,r)$: there are $2^{m-r}$ finite ramification points, all wild, with index and different exponent $2^{r-2}$, all over the single finite branch value $1$. The $0/1$ lift of $P_{m,r}$ satisfies an identity in $\Z[X]$ that evaluates it over every Galois ring $\GR(2^k,m)$ with $m-1$ squarings, one general multiplication and $r+5$ additions or subtractions. For every quadratic optimal APN permutation $p$ of $\F_q$ (in particular $P_{m,r}$), every $k\ge2$ and every coefficient lift $F$ of $p$ to $\GR(2^k,m)$, we determine the differential rows of $F$ exactly in the unit directions whose residue lies outside the kernel $\{0,\rho\}$ of the linear part of $p'$ (they are the residue rows of $p$; $\rho=1$ for $P_{m,r}$) and in the directions in $2^{k-1}\GR(2^k,m)\setminus\{0\}$ (every nonzero entry equals $2q^{k-1}$). We also give a self-contained proof of the Li--Wang criterion for a Gold function plus a linearised polynomial, which explains why the inner map of $P_{m,r}$ is not a monomial, and record a two-word, two-layer transition that a simple linear layer forbids at the residue level.
\end{abstract}
\begin{keyword}
APN functions \sep permutation polynomials \sep Galois rings \sep ramification \sep Gold functions \sep differential uniformity
\MSC[2020] 11T06 \sep 11T71 \sep 94A60 \sep 14H30
\end{keyword}
\end{frontmatter}


\section{Introduction}
R{\o}njom and Sandrib \cite{RS} asked for polynomial S-boxes over Galois rings that are APN on the residue field, have the smallest possible differential uniformity over the ring, and can be reused across extension degrees. Bartoli and St\u{a}nic\u{a} \cite{BS} reduced the lifting question to the formal derivative of the reduced representative on $\F_q$ and conjectured that it always has a zero. Explicit counterexamples are given in a public code repository by Stone \cite{Stone}, and \cite{TIT} shows that every quadratic APN permutation class in odd dimension $m\ge5$ contains a representative whose derivative has no zero (an \emph{optimal} representative). That existence theorem uses coefficients in $\F_q$. The present paper does not use it.

We study one explicit family with coefficients in $\F_2$. It is a uniform, parameter-dependent construction: one rule gives $P_{m,r}\in\F_2[X]$ for every admissible pair $(m,r)$, but the exponents depend on $m$, and no fixed polynomial is claimed to work for all $m$ (for instance $P_{5,3}$ does not permute $\F_{2^7}$). For a fixed pair, the statements on coefficient lifts hold over $\GR(2^k,m)$ for every $k\ge2$. We record the following.
\begin{enumerate}
\item[(A)] The family $P_{m,r}$ and its derivative (Theorem~\ref{thm:A}).
\item[(B)] Its ramification over the algebraic closure, for all admissible parameters: critical points, ramification indices, different exponents, branch values (Theorem~\ref{thm:B}).
\item[(C)] An identity in $\Z[X]$ for the $0/1$ lift that evaluates it over $\GR(2^k,m)$ with one general multiplication (Proposition~\ref{prop:C}).
\item[(R)] For every coefficient lift of a quadratic optimal APN permutation, in particular of $P_{m,r}$, exact differential rows in two classes of directions (Theorem~\ref{thm:R}).
\item[(W)] A two-word statement about one simple linear layer (Proposition~\ref{prop:W}).
\end{enumerate}

\paragraph{What is and is not new.} For $m=5$, $P_{5,3}$ is Stone's hand example $(x+x^8+x^{16})^3$ \cite{Stone}. His further examples, for $m=5,7,9$, are given as value tables and are stated there to be relabellings of $x^{11}$, $x^{11}$ and $x^{5}$; they are not of the form $P_{m,r}$, which is linearly equivalent to $x^3$ (for $m=5,7$ the algebraic degree differs, and for $m=9$ the functions $x^3$ and $x^5$ lie in different classes \cite[Table~1]{YKBL}). Optimal representatives with $\F_q$ coefficients exist in every quadratic APN permutation class \cite{TIT}, and the closed forms $P_{e,\theta}$ of \cite[App.~A]{TIT} with $e=1$ also have a single finite branch value (Remark~\ref{rem:tit}); for $e\ge2$ they have $2^{e-1}$. So neither existence nor a single branch value alone is claimed as new. What we add is the combination: a single rule giving coefficients in $\F_2$ and $3r-2$ terms for every admissible $(m,r)$; the full list of ramification data (number of critical points, local indices, different exponents) for all these parameters; and the integer product form of the natural lift. The permutation criterion for a Gold function plus a linearised polynomial (Lemma~\ref{lem:LW}) is due to Li and Wang \cite[Thm.~4]{LW}; we give a self-contained proof for every Gold exponent and use it to explain why the inner map of $P_{m,r}$ cannot be a monomial (Corollary~\ref{cor:P}). The ring statements of Section~\ref{sec:ring} are proved directly and hold for every quadratic optimal APN permutation; for the remaining directions we quote upper bounds from \cite{RS}, which yields the ring corollary of \cite{TIT}.

\section{Preliminaries}
Throughout, $m\ge5$ is odd, $q=2^m$, and $\sigma(x)=x^2$. For $f\in\F_q[X]$, $\red(f)$ is its reduction modulo $X^q-X$, and $[X^e]f$ denotes the coefficient of $X^e$ in $f$.

A map $f:\F_q\to\F_q$ is \emph{APN} if for every $a\neq0$ and every $b$ the equation $f(x+a)+f(x)=b$ has at most two solutions. A \emph{linearised polynomial} over $\F_q$ is $M=\sum_{i<m}\mu_iX^{2^i}$ with $\mu_i\in\F_q$; the associated map is additive, that is $\F_2$-linear, and every $\F_2$-linear map of $\F_q$ arises in this way. (Only the maps $x\mapsto ax$ are $\F_q$-linear.) The $\F_2$-linear maps commuting with $\sigma$ are the polynomials $\kappa(\sigma)=\sum\kappa_i\sigma^i$ with $\kappa\in\F_2[t]/(t^m-1)$, i.e.\ the linearised polynomials with coefficients in $\F_2$; $\kappa(\sigma)$ is bijective iff $\gcd(\kappa,t^m-1)=1$.

A \emph{Dembowski--Ostrom (DO) polynomial} is an $\F_q$-linear combination of monomials $X^{2^i+2^j}$, $0\le i<j<m$. A \emph{quadratic} polynomial is a DO polynomial plus a linearised polynomial plus a constant. The derivative of $X^{2^i+2^j}$ is $X^{2^j}$ if $i=0$ and $0$ otherwise, and that of $X^{2^i}$ is $1$ if $i=0$ and $0$ otherwise. Hence for quadratic $p$,
\begin{equation}\label{eq:der}
p'=\epsilon+K(X),\qquad \epsilon=[X]p\in\F_q,\quad K=\sum_{j\ge1}([X^{1+2^j}]p)\,X^{2^j}\ \text{linearised}.
\end{equation}

\begin{definition}\label{def:opt}
A reduced $p\in\F_q[X]$ is \emph{derivative-optimal} if $p'$ has no zero on $\F_q$ and every nonempty fibre of $x\mapsto p'(x)$ on $\F_q$ has two elements. An \emph{optimal APN permutation} is an APN permutation of $\F_q$ whose reduced polynomial is derivative-optimal.
\end{definition}
For quadratic $p$, \eqref{eq:der} shows that $p$ is derivative-optimal iff $\ker K=\{0,\rho\}$ for some $\rho\ne0$ and $\epsilon\notin K(\F_q)$; the fibres of $p'$ are the cosets $\{x,x+\rho\}$.

\begin{lemma}\label{lem:O}
Let $p$ be a reduced quadratic polynomial with coefficients in $\F_2$, and let $\kappa\in\F_2[t]$ be the polynomial $\sum_{j\ge1}([X^{1+2^j}]p)\,t^j$, so that $K=\kappa(\sigma)$. Then $p$ is derivative-optimal iff $\epsilon=1$ and $\gcd(\kappa,t^m-1)=t+1$.
\end{lemma}
\begin{proof}
Fix a normal element $\beta$; then $f\mapsto f(\sigma)\beta$ identifies $\F_2[t]/(t^m-1)$ with $\F_q$, and $\kappa(\sigma)$ becomes multiplication by $\kappa$. Its kernel has dimension $\deg\gcd(\kappa,t^m-1)$ and its image is the ideal generated by this gcd. The fibres of $p'$ have size two iff the gcd has degree one, i.e.\ equals $t+1$. If $\epsilon=0$, $0$ is a value. If $\epsilon=1$: $\Tr(\beta)=1$ (the conjugates of a normal element are linearly independent, so their sum is not $0$), so $1=\sum_i\sigma^i\beta$ corresponds to $1+t+\dots+t^{m-1}$, which is not in $(t+1)$ since its value at $1$ is $m\equiv1$.
\end{proof}

\section{The family}\label{sec:A}
Call an odd $r$ with $3\le r<m$ and $\gcd(r,m)=1$ \emph{admissible} ($r=m-2$ always is). Put $T_r=\sum_{i<r}X^{2^i}$ and
\begin{gather*}
V=\sum_{i=0}^{r-3}X^{2^i},\quad W=X^{2^{r-2}},\quad Y=X^{2^{m-2}},\quad Z=X^{2^{m-1}},\\
D=1+W+Y,\quad E=1+V+Y+Z.
\end{gather*}
\begin{theorem}\label{thm:A}
Let $r$ be admissible and $P_{m,r}:=\red\bigl((T_r^3)^{2^{m-2}}\bigr)$; we write $P=P_{m,r}$ when $(m,r)$ is fixed.
\begin{enumerate}
\item $P_{m,r}=V+W+WV+WY+WZ+YV+YZ$; it has coefficients in $\F_2$, exactly $3r-2$ terms and degree $3q/4$.
\item As a map of $\F_q$, $P_{m,r}=\sigma^{m-2}\circ X^3\circ T_r$; it is an APN permutation.
\item $P'_{m,r}=D=1+X^{2^{r-2}}+X^{2^{m-2}}$; $P_{m,r}$ is derivative-optimal, with fibres $\{x,x+1\}$ and image of $P'_{m,r}$ equal to $1+\ker\Tr$.
\item $P_{m,r}+1=D\cdot E$ in $\F_2[X]$, and $E'=1$.
\end{enumerate}
\end{theorem}
\begin{proof}
(2) $T_r$ is $\tau(\sigma)$ with $\tau=(t^r-1)/(t-1)$. Since $\gcd(t^r-1,t^m-1)=t-1$ and $\tau(1)=r\equiv1$, $\tau$ is a unit and $T_r$ is bijective. $X^3$ is an APN permutation of $\F_q$ for odd $m$, and $\sigma^{m-2}$ is additive and bijective; composing with additive bijections on both sides preserves both properties.

(1) $T_r^3=T_r^2T_r=\sum_{i,j<r}X^{2^{i+1}+2^j}$. The diagonal pairs $i+1=j$ give the linear terms $X^{2\cdot2^j}=X^{2^{j+1}}$ ($1\le j\le r-1$). A two-element set $\{a,b\}\subseteq[1,r-1]$ arises from $(i+1,j)=(a,b)$ and from $(b,a)$ and cancels; the survivors are $\{0,b\}$ ($1\le b\le r$) and $\{a,r\}$ ($1\le a\le r-1$); the set $\{0,r\}$ occurs once, from $(i+1,j)=(r,0)$. So $T_r^3$ has $r-1$ linear and $2r-1$ quadratic terms, all of degree $<q$. Raising to the power $2^{m-2}$ and reducing rotates every index by $-2$ modulo $m$, injectively. The result has the linear terms $X^{2^i}$, $0\le i\le r-2$ ($=V+W$), and the quadratic terms with index sets $\{m-2,b-2\}$, $1\le b\le r$, and $\{a-2,r-2\}$, $1\le a\le r-1$, which are the supports of $YZ$, $YV$, $WY$ and of $WZ$, $WV$. The count is $3r-2$; the largest exponent is $2^{m-1}+2^{m-2}=3q/4$.

(3) The odd exponents are $1$ (in $V$), $2^{r-2}+1$ (in $WV$) and $2^{m-2}+1$ (in $YV$), so $P'=D$. Here $\kappa=t^{r-2}+t^{m-2}=t^{r-2}(1+t^{m-r})$ and $\gcd(\kappa,t^m-1)=t^{\gcd(m-r,m)}-1=t+1$; Lemma~\ref{lem:O} applies. Directly, $x^{2^{r-2}}+x^{2^{m-2}}=(x^{2^r}+x)^{2^{m-2}}$ has kernel $\F_2$ and image $\ker\Tr$.

(4) In $\F_2[X]$, $DE=1+V+W+Z+WV+WY+WZ+YV+Y^2+YZ$ (the two terms $Y$ cancel), and $Y^2=Z$; hence $DE=1+P_{m,r}$. $E'=V'=1$.
\end{proof}

\subsection*{Why the inner map is not a monomial}
\begin{lemma}[Li--Wang {\cite[Thm.~4]{LW}}]\label{lem:LW}
Let $1\le h<m$ with $\gcd(h,m)=1$ and let $M=\sum_{i<m}\mu_iX^{2^i}$ be a reduced linearised polynomial over $\F_q$. Then $X^{2^h+1}+M$ permutes $\F_q$ iff $M=\gamma^{2^h}X+\gamma X^{2^h}$ for some $\gamma\in\F_q$, i.e.\ $X^{2^h+1}+M=(X+\gamma)^{2^h+1}+\gamma^{2^h+1}$.
\end{lemma}
Li and Wang state this for nonzero $M$, i.e.\ $\gamma\neq0$; $\gamma=0$ is the Gold function itself, a permutation since $m$ is odd. They write out the proof for $h=1$ \cite[Thm.~2]{LW}, where the binary weights of the relevant exponents already separate the cyclotomic classes, and note that it extends to $\gcd(h,m)=1$. For general $h$ two classes can have the same weight, so we give a short self-contained proof that includes this comparison.
\begin{proof}
For $a\ne0$ and $z=x/a$, $(x+a)^{2^h+1}+x^{2^h+1}+M(a)=a^{2^h+1}(z^{2^h}+z+1)+M(a)$, and $z^{2^h}+z$ ranges over $\ker\Tr$ (its kernel is $\F_{2^{\gcd(h,m)}}=\F_2$; this is the only use of $\gcd(h,m)=1$). As $\Tr(1)=1$ ($m$ odd), $X^{2^h+1}+M$ is a permutation iff $\Tr\bigl(M(a)a^{-(2^h+1)}\bigr)=0$ for all $a\ne0$. With $c_i=2^i-2^h-1$ this reads $\sum_i\Tr(\mu_ia^{c_i})=0$ on $\F_q^*$. Modulo $Q=q-1$, multiplication by $2$ rotates $m$-bit words (bit $j$ is the coefficient of $2^j$); the orbits of $c\mapsto2c$ on $\Z/Q$ are the cyclotomic classes. The word of $c_h=-1$ has its only zero at position $0$, that of $c_0=-2^h$ at position $h$, so $c_0=2^hc_h$. For $0<i<h$ the ones of $c_i$ occupy $[1,i-1]\cup[h,m-1]$ (weight $m-h+i-1$), for $h<i<m$ they occupy $[0,h-1]\cup[h+1,i-1]$ (weight $i-1$); both weights are at most $m-2$. Every such word contains both digits and has one or two cyclic runs of zeros. A word fixed by a nontrivial rotation consists of $m/d$ copies of a block, $d$ a proper divisor of the odd number $m$, so its number of zero runs would be a positive multiple of $m/d\ge3$; hence every class has $m$ elements, and none is the class of $0$. Words of different weight lie in different classes, so it remains to compare $c_u$ and $c_v$ with $0<u<h<v=m-h+u$. Their zero sets are $\{0\}\cup[u,h-1]$ and $\{h\}\cup[v,m-1]$; the first is a single run iff $u=1$, the second iff $v=h+1$, and both happen only if $m=2h$. If exactly one of them is a single run, the two words have different numbers of zero runs and lie in different classes. Otherwise ($u\ge2$, $v\ge h+2$) both have zero runs of lengths $1$ and $h-u$. If $h-u\ge2$, a rotation matches the isolated zeros, and the runs of ones that follow them have lengths $u-1$ and $v-h-1$, so $v=u+h$ and $m=2h$. If $h-u=1$, the gaps between the two zeros are $\{h-2,m-h\}$ and $\{h,m-h-2\}$, equal only if $m=2h$. As $m$ is odd, the classes of the $c_i$, $i\ne0,h$, are pairwise distinct and distinct from that of $c_h$. So the left side is a polynomial in $a$ with exponents in $[1,Q-1]$ whose coefficients are $\mu_i^{2^j}$ at $a^{c_i2^j}$, except on the class of $c_h$, where they are $\mu_h^{2^j}+\mu_0^{2^{j-h}}$. A polynomial of degree $<Q$ vanishing on $\F_q^*$ is zero; hence $\mu_i=0$ for $i\notin\{0,h\}$ and $\mu_0=\mu_h^{2^h}$, i.e.\ $M=\gamma^{2^h}X+\gamma X^{2^h}$ with $\gamma=\mu_h$. Conversely, for such $M$ the identity in the statement and the bijectivity of $X^{2^h+1}$ ($m$ odd) show that $X^{2^h+1}+M$ permutes $\F_q$.
\end{proof}

\begin{corollary}\label{cor:P}
Let $G=X^{2^h+1}$ with $\gcd(h,m)=1$, let $N$ be an $\F_2$-linear bijection and $L$ an $\F_2$-linear map of $\F_q$, $b\in\F_q^*$ and $0\le j<m$. No permutation of the form $N\circ G\circ(bX^{2^j})+L+\text{const}$ is derivative-optimal. Hence an optimal APN permutation in the extended-affine class of $G$, written $N\circ G\circ B+L+\text{const}$ with $B$ an $\F_2$-linear bijection (an inner translation only adds an affine map, as $G$ is quadratic), has an inner map $B$ that is not a monomial.
\end{corollary}
\begin{proof}
$G\circ(bX^{2^j})=b^{2^h+1}\sigma^j\circ G$, and $N\circ(b^{2^h+1}\sigma^j)$ is again an $\F_2$-linear bijection, so it suffices to treat $N\circ G+L$. If it permutes $\F_q$, so does $G+N^{-1}L$, and Lemma~\ref{lem:LW} gives $N\circ G+L=H(X+\gamma)+\text{const}$ with $H=\red(N\circ G)$. As $\deg H<q$, $H(X+\gamma)$ is reduced, and its formal derivative is $H'(X+\gamma)$. Every exponent of $H$ has binary weight two, so $H$ has no term $X$ and $H'(0)=0$ by \eqref{eq:der}; the derivative $H'(X+\gamma)$ of the permutation vanishes at $\gamma$.
\end{proof}
For $m\ge5$ and representations without affine part the inner map is determined up to a monomial: replacing $h$ by $m-h$ changes $G$ only by a power of $\sigma$ on the output, so we may assume $h<m/2$, and if $N_1\circ G\circ B_1=N_2\circ G\circ B_2$, then $B_1B_2^{-1}$ is a monomial $aX^{2^u}$ \cite[Thm.~4.1]{KZ}. With affine parts the corresponding statement concerns the extended-affine automorphism group of $G$; we do not use it. In Theorem~\ref{thm:A} the inner map is $T_r$.

\section{Ramification}\label{sec:B}
Let $s=m-r$ (even), $\ell=2^{r-2}$ and $g=X^{2^s}+X+1\in\F_2[X]$. We regard $P=P_{m,r}$ as a map $\mathbb P^1\to\mathbb P^1$ over $\overline{\F}_2$ of degree $n=3q/4$; it is separable, as $P'\ne0$. At a finite point $\alpha$, the ramification index $e_\alpha$ is the multiplicity of $\alpha$ as a root of $P-P(\alpha)$, and the different exponent is $d_\alpha=\operatorname{ord}_\alpha P'$; in characteristic two $d_\alpha\ge e_\alpha$ when $e_\alpha$ is even (wild ramification), so $d_\alpha$ is not $e_\alpha-1$ in general. A \emph{ramification point} is a point with $e_\alpha>1$; its image is a \emph{branch value}.
\begin{theorem}\label{thm:B}
\begin{enumerate}
\item $P_{m,r}'=g^{\ell}$ and $P_{m,r}+1=g^{\ell}E$.
\item $g$ is separable; its $2^s$ roots lie in $\F_{2^{2s}}\setminus(\F_{2^s}\cup\F_q)$.
\item $\gcd(g,E)=1$.
\end{enumerate}
Consequently the finite ramification points are the $2^s$ roots of $g$; each has $e_\alpha=\ell$ and $d_\alpha=\ell$ and maps to $1$. The point $\infty$ has $e_\infty=n$ and $d_\infty=2n-2-\deg P'=5q/4-2$. The branch values are exactly $1$ and $\infty$. The fibre over $1$ consists of the roots of $g$ and $2^{m-1}$ simple points, exactly one of which is $\F_q$-rational.
\end{theorem}
\begin{proof}
(1) $g^{\ell}=X^{2^{m-2}}+X^{2^{r-2}}+1=D$; use Theorem~\ref{thm:A}(3),(4). (2) $g'=1$. If $g(\alpha)=0$ then $\alpha^{2^s}=\alpha+1$, so $\alpha\notin\F_{2^s}$ and $\alpha^{2^{2s}}=\alpha$. As $\gcd(2s,m)=1$, $\F_{2^{2s}}\cap\F_q=\F_2$, and $g(0)=g(1)=1$. (3) Modulo $g$, $X^{2^{s+j}}\equiv X^{2^j}+1$, so $Y\equiv X^{2^{r-2}}+1$, $Z\equiv X^{2^{r-1}}+1$ and $E\equiv1+T_r$. If $g(\alpha)=T_r(\alpha)+1=0$, then $T_r^2+T_r=X+X^{2^r}$ gives $\alpha^{2^r}=\alpha$, so $\alpha\in\F_{2^{\gcd(r,2s)}}=\F_2$ ($r$ odd, $\gcd(r,s)=\gcd(r,m)=1$), contradicting $g(0)=g(1)=1$.

$P'$ vanishes only at the roots of $g$, so these are the only finite ramification points. At a root $\alpha$, $P(\alpha)=1$ by (1), and $\alpha$ is a root of multiplicity exactly $\ell$ of $P+1=g^{\ell}E$, since $g$ is separable and $E(\alpha)\neq0$ by (3); so $e_\alpha=\ell$, and $d_\alpha=\operatorname{ord}_\alpha g^{\ell}=\ell$. Riemann--Hurwitz for the separable degree-$n$ map of genus zero gives $\sum_\alpha d_\alpha+d_\infty=2n-2$; the finite sum is $2^s\ell=q/4=\deg P'$, so $d_\infty=3q/2-2-q/4=5q/4-2$, and $e_\infty=n$ since $P$ is a polynomial. $E$ has degree $2^{m-1}$, is separable ($E'=1$), and $P'=g^\ell$ does not vanish at its roots by (3), so they are simple points of the fibre over $1$; the degrees match, $2^s\ell+2^{m-1}=q/4+q/2=n$. As $P$ permutes $\F_q$, $P(x)=1$ has exactly one rational solution, which is not a root of $g$ by (2).
\end{proof}
All ramification is wild ($\ell$ and $n$ are even), and at the finite points $d_\alpha=e_\alpha$ is the least value allowed for wild ramification. Galois closures, monodromy and higher ramification groups are not determined here. For $r=m-2$ the finite ramification points are the four roots of $X^4+X+1$ for every $m$.

\begin{remark}[comparison with \cite{TIT}]\label{rem:tit}
Let $1\le e\le(m-1)/2$, $\theta\in\F_q^\times$ with $\Tr(\theta)=0$, $L_\theta=\sum_i\theta^{2^i}X^{2^i}$ and $P_{e,\theta}=X^{2^{e+1}+2}+(X^{2^{e+1}}+X^2+1)L_\theta$. If $\gcd(e,m)=1$, $P_{e,\theta}$ is, by \cite[App.~A]{TIT}, an optimal APN permutation linearly equivalent to $X^{2^e+1}$. Without this condition it need not be APN or derivative-optimal: for $m=9$, $e=3$ and $\theta$ a root of the irreducible polynomial $X^9+X^4+1$ (its trace is $0$), $P_{e,\theta}$ permutes $\F_q$, but its differential uniformity is $8$ and the nonempty fibres of $P_{e,\theta}'$ have $8$ elements. The count of branch values below uses neither condition. Put $g_e=X^{2^e}+X+1$, so $g_e^2=X^{2^{e+1}}+X^2+1$ and $P_{e,\theta}'=\theta\,g_e^2$. At a root $\alpha$ of $g_e$, $P_{e,\theta}(\alpha)=\alpha^{2^{e+1}+2}=(\alpha^2+\alpha)^2$, using $\alpha^{2^e}=\alpha+1$. The roots form a coset $\alpha_0+\F_{2^e}$, and $u\mapsto(\alpha_0+u)^2+(\alpha_0+u)=(\alpha_0^2+\alpha_0)+(u^2+u)$ is two-to-one onto a coset of the trace-zero hyperplane of $\F_{2^e}$; squaring is injective. Hence the finite branch values of $P_{e,\theta}$ are exactly $2^{e-1}$ in number. For $e=1$ there is one, and explicitly
\[
P_{1,\theta}+1=(X^2+X+1)^2\bigl((X+1)^2+L_\theta\bigr).
\]
So a single finite branch value is not special to Theorem~\ref{thm:A}: the maps $P_{1,\theta}$, which lie in the class of $X^3$ like $P_{m,r}$, have it, with coefficients outside $\F_2$ ($\Tr(\theta)=0$ excludes $\theta=1$) and two finite critical points. Theorem~\ref{thm:A} has coefficients in $\F_2$ and $2^{m-r}$ critical points of index $2^{r-2}$.
\end{remark}

\section{The integer product form}\label{sec:C}
\begin{proposition}\label{prop:C}
Let $\hat P\in\Z[X]$ be the $0/1$ lift of $P_{m,r}$ and $\hat V,\hat W,\hat Y,\hat Z,\hat D,\hat E$ the $0/1$ lifts of $V,\dots,E$. Then
\[
\hat P=\hat D\hat E-1-2(\hat Y+\hat Z)\quad\text{in }\Z[X].
\]
Over every Galois ring $\GR(2^k,m)$, $k\ge1$, $\hat P(x)$ can be evaluated with $m-1$ squarings, one general multiplication and $r+5$ additions or subtractions.
\end{proposition}
\begin{proof}
In $\Z[X]$, $\hat D\hat E=1+\hat V+\hat W+2\hat Y+\hat Z+\hat W\hat V+\hat W\hat Y+\hat W\hat Z+\hat Y\hat V+\hat Y^2+\hat Y\hat Z$ and $\hat Y^2=\hat Z$. Hence $\hat D\hat E-1-2(\hat Y+\hat Z)=\hat V+\hat W+\hat W\hat V+\hat W\hat Y+\hat W\hat Z+\hat Y\hat V+\hat Y\hat Z$; its monomials are pairwise distinct with coefficient $1$, and by Theorem~\ref{thm:A}(1) this is $\hat P$. Schedule: $x^{2^i}$ for $1\le i\le m-1$ by squaring; $S=\hat Y+\hat Z$ (one addition), $\hat D=1+\hat W+\hat Y$ (two), $\hat E=(1+\hat V)+S$ ($r-1$), $\hat D\hat E$ (one multiplication), $\hat D\hat E-(1+S+S)$ (three). Squarings in the ring are counted as operations, not as free.
\end{proof}
This is an upper bound for the natural lift. We do not claim that it is minimal, and we say nothing about the cost of a complete cipher.

\section{Coefficient lifts}\label{sec:ring}
Let $R=\GR(2^k,m)$, $k\ge2$, with residue map $x\mapsto\bar x$ onto $\F_q$. For $p\in\F_q[X]$, a \emph{coefficient lift} is any $F\in R[X]$ with $\bar F=p$. Differences are taken in $R$: for $\alpha,\beta\in R$, $\delta_F(\alpha,\beta)=\#\{x\in R:F(x+\alpha)-F(x)=\beta\}$, and $\Delta_F=\max_{\alpha\ne0,\beta}\delta_F(\alpha,\beta)$.

\begin{proposition}\label{prop:lift}
If $p$ permutes $\F_q$ and $p'$ has no zero on $\F_q$, every coefficient lift of $p$ permutes $R$.
\end{proposition}
\begin{proof}
If $F(x)=F(y)$, then $p(\bar x)=p(\bar y)$, so $y=x+\eta$ with $\eta\in2R$. If $\eta\neq0$, let $\eta\in2^jR\setminus2^{j+1}R$, $j\ge1$. Then $F(x+\eta)-F(x)=\eta F'(x)+\eta^2H_0$ with $H_0\in R$; $F'(x)$ is a unit since $\overline{F'(x)}=p'(\bar x)\ne0$, so $\eta F'(x)\notin2^{j+1}R$, while $\eta^2H_0\in2^{2j}R\subseteq2^{j+1}R$. Hence $F(x+\eta)\ne F(x)$.
\end{proof}
This is the sufficiency half of \cite[Thm.~2]{RS}, included for completeness.

\begin{theorem}\label{thm:R}
Let $p$ be a quadratic optimal APN permutation of $\F_q$, with $p'=\epsilon+K$ as in \eqref{eq:der} and $\ker K=\{0,\rho\}$, and let $F$ be any coefficient lift.
\begin{enumerate}
\item (Unit directions off $\rho$.) If $\bar\alpha\notin\{0,\rho\}$, then $\delta_F(\alpha,\beta)=\#\{y\in\F_q:p(y+\bar\alpha)+p(y)=\bar\beta\}\le2$ for every $\beta\in R$. There are $(q-2)q^{k-1}$ such $\alpha$.
\item (Top ideal.) If $\alpha\in2^{k-1}R\setminus\{0\}$, then every nonzero entry $\delta_F(\alpha,\beta)$ equals $2q^{k-1}$.
\end{enumerate}
\end{theorem}
\begin{proof}
(1) Put $\mathcal D(X)=F(X+\alpha)-F(X)-\beta$. Its reduction is $p(X+\bar\alpha)+p(X)+\bar\beta$, whose derivative $p'(X+\bar\alpha)+p'(X)=K(\bar\alpha)$ is a nonzero constant. So every root of the reduction in $\F_q$ is simple and lifts to exactly one root of $\mathcal D$ in $R$ (Hensel), and every root of $\mathcal D$ reduces to a root. The bound $2$ is the APN property of $p$.
(2) Write $\alpha=2^{k-1}u$, $u$ a unit. By Taylor expansion with integer binomial coefficients, $F(x+\alpha)-F(x)=\alpha F'(x)+\alpha^2(\dots)=\alpha F'(x)$, as $\alpha^2=0$ for $k\ge2$. Since $\alpha\cdot2R=0$, $\alpha F'(x)$ depends only on $\bar x$ and corresponds to $\bar u\,p'(\bar x)$ under $2^{k-1}R\cong\F_q$ (multiplication by $2^{k-1}$ induces an additive isomorphism $R/2R\to2^{k-1}R$, as the annihilator of $2^{k-1}$ in $R$ is $2R$). The map $\bar x\mapsto\bar u\,p'(\bar x)$ is two-to-one onto its image, and each residue class has $q^{k-1}$ elements.
\end{proof}
Part~(2) of the theorem uses only that $p$ is derivative-optimal; part~(1) uses that $p$ is quadratic and APN with $\ker K=\{0,\rho\}$. The theorem gives exact rows in these two classes of directions. In the remaining directions only upper bounds are used, from \cite{RS}: if $\bar\alpha=\rho$, \cite[Prop.~4(1)]{RS} gives $\delta_F(\alpha,\beta)\le2q^{k-1}$; if $\alpha\in2R\setminus2^{k-1}R$, \cite[Prop.~5]{RS}, applied with the fibres of $p'$ of size two, gives the same bound. Since by part~(2) some entry equals $2q^{k-1}$, $\Delta_F=2q^{k-1}$ for every coefficient lift; this is the ring corollary of \cite{TIT} (the lower bound also follows from \cite[Thm.~3]{RS}). For $m=5$ we computed the full tables of $\hat P_{5,3}$ and of one random coefficient lift of $P_{5,3}$ (drawn with a fixed seed; see the accompanying files) for $k=2,3$: all rows agree with Theorem~\ref{thm:R}, and the largest entry is $2q^{k-1}$ ($64$ and $2048$); in the directions over $\rho$ the row maxima of $\hat P_{5,3}$ are $64$ ($k=2$) and $128$ ($k=3$).

\section{Two words, two layers}\label{sec:W}
Let $p$ be as in Theorem~\ref{thm:R}. Replacing $p$ by $c\,p(\rho x)$, $c\ne0$, we may assume $\rho=1$ and $\{p(y+1)+p(y):y\in\F_q\}=\{z:\Tr(z)=1\}$: the set $\{p(y+\rho)+p(y)\}$ is an affine hyperplane (a coset of an $\F_2$-subspace of dimension $m-1$, since $p$ is quadratic and APN) not containing $0$ (since $p$ is a permutation), hence of the form $\{z:\Tr(vz)=1\}$, and $c=v$ works.
\begin{proposition}\label{prop:W}
With this normalisation, let $a,b\in\F_q$ be distinct and nonzero with $\Tr(a)=\Tr(b)=0$, $\mathbf B=\begin{psmallmatrix}a&b\\b&a+b\end{psmallmatrix}$ and $\mathbf M=\mathbf B^{-1}$. In the two-word map $x\mapsto\mathbf M\cdot(p(x_1),p(x_2))$, there is no pair of consecutive layers in which the S-box input differences $\alpha$ and then $\mathbf M\beta$ both lie in $\{0,1\}^2\setminus\{0\}$; here $\beta$ is any output difference of $\alpha$, i.e.\ $p(y_i+\alpha_i)+p(y_i)=\beta_i$ has a solution for $i=1,2$.
\end{proposition}
\begin{proof}
$\det\mathbf B=a^2+ab+b^2\ne0$ because $X^2+X+1$ has no root in $\F_q$ for odd $m$. If $\mathbf M\beta\in\{0,1\}^2\setminus\{0\}$, then $\beta=\mathbf B\gamma$ with $\gamma\in\{0,1\}^2\setminus\{0\}$, so $\beta\in\{(a,b),(b,a+b),(a+b,a)\}$: both entries nonzero with trace $0$. Word by word, $\alpha_i=0$ forces $\beta_i=0$, and $\alpha_i=1$ forces $\Tr\beta_i=1$; both are excluded.
\end{proof}
The statement is algebraic: it concerns residue-level differences in two words and two consecutive layers. Reading products of transition probabilities as characteristic probabilities requires independent uniform round keys, which we do not otherwise use. We give no bound on differentials, hulls or fixed-key behaviour, and directions in $2R$ (residue difference zero) are not covered.

\section{Relation to prior work}
\begin{itemize}
\item \cite{BS}: the reduced lifting and critical-point conjectures; Theorem~\ref{thm:A} gives explicit counterexamples for every odd $m\ge5$. \cite{Stone} gives a hand example ($P_{5,3}$) and value tables for $m=5,7,9$ in other classes; \cite{TIT} gives counterexamples in every class and every odd dimension $m\ge5$, with coefficients in $\F_q$.
\item \cite{RS}: the lifting criterion (Thm.~2; Proposition~\ref{prop:lift} is its sufficiency half), the lower bound (Thm.~3; for the lifts considered here it also follows from Theorem~\ref{thm:R}(2)) and the upper bounds used after Theorem~\ref{thm:R} (Props.~4 and~5). Their Conjecture~2 (an APN permutation of $\F_{2^m}$ has no lift that permutes $\GR(2^k,m)$) is refuted by Theorem~\ref{thm:A}(2),(3) together with Proposition~\ref{prop:lift}, for every odd $m\ge5$ and every $k\ge2$; it is also refuted by the examples of \cite{Stone} for $m=5,7,9$ and by \cite{TIT} for every odd $m\ge5$.
\item \cite[App.~A]{TIT}: closed forms $P_{e,\theta}$ with $\F_q$ coefficients in every Gold class; Remark~\ref{rem:tit} compares their ramification with Theorem~\ref{thm:B}; the case $e=1$ also has one finite branch value.
\item \cite[Thm.~4]{LW}: Lemma~\ref{lem:LW} (proof written out there for $h=1$, Thm.~2). \cite[Thm.~9]{Chen} gives a trace criterion for Gold-type polynomials plus linearised polynomials over finite fields of characteristic two; for the base field $\F_2$, i.e.\ for $X^{2^h+1}+M$ over $\F_{2^m}$, the criterion of Lemma~\ref{lem:LW} is the earlier result of \cite{LW}.
\item \cite{KZ}: linear equivalences $G\circ L=N\circ G$ between Gold functions are monomial for $m\ge5$ (Thm.~4.1 there); used after Corollary~\ref{cor:P} for representations without affine part.
\item \cite{BL,YKBL}: searches and classifications of quadratic APN functions; the permutations of Theorem~\ref{thm:A} lie in the Gold class of $X^3$.
\end{itemize}

\section*{Methods and reproducibility}
The script \texttt{run\_quick.sh} (in the folder that accompanies this paper) recomputes every finite computation quoted above and calibrates the general identities on finite parameter ranges (for Theorem~\ref{thm:A} and Proposition~\ref{prop:C} all admissible $(m,r)$ with $m\le101$; the other ranges are listed in \texttt{REPRODUCE.md}), from the files in that folder alone (Python~3 with NumPy). These checks do not replace the proofs, which are given in the text. Degree-minimisation results about related polynomials are kept in a separate research note and are not used here.

\emph{Use of AI tools in the research.} Claude (Anthropic), including the agentic tool Claude Code, and ChatGPT (OpenAI) were used in proof exploration, in drafting and revising arguments, in writing and running the finite verification scripts, and (Claude) in writing the Lean formalization described below; ChatGPT was also used for repeated adversarial reviews of the draft. The general results rest on the proofs given in the text; the programs provide finite calibrations, not replacements for those proofs. The checks and the scope of the independently written implementations are documented in the accompanying files. These checks and reviews are internal; they are not external peer review.

\emph{Formalization.} The following statements are formalized in Lean~4 with Mathlib; the formal proofs use only the axioms \texttt{propext}, \texttt{Classical.choice} and \texttt{Quot.sound}, and contain no \texttt{sorry}. Lemma~\ref{lem:O}, for every polynomial over $\F_2$ whose exponents are $0$, $2^i$ or $2^i+2^j$ with $i<j<m$ (the formal proof replaces the normal basis by a count of roots: the kernel of $\phi(\sigma)$ has $2^{\deg\phi}$ elements for every divisor $\phi$ of $t^m-1$); Theorem~\ref{thm:A}, all four items, with $\Tr$ written as $\sum_{i<m}y^{2^i}$; the identity of Proposition~\ref{prop:C} in $\Z[X]$ and in every commutative ring; items (1)--(3) of Theorem~\ref{thm:B} with the root multiplicities $e_\alpha=d_\alpha=\ell$ and $P(\alpha)=1$; Proposition~\ref{prop:lift} and Theorem~\ref{thm:R} for every finite commutative ring $R$ with a surjection $\pi$ onto $\F_q$ such that elements outside $\ker\pi$ are units, and, in part~(2), for every $\alpha\in\ker\pi$ whose annihilator is $\ker\pi$ (the nonzero elements of $2^{k-1}\GR(2^k,m)$ when $k\ge2$); and Proposition~\ref{prop:W}, with the normalisation made before it taken as a hypothesis. Not formalized: Lemma~\ref{lem:LW}, Corollary~\ref{cor:P}, Remark~\ref{rem:tit}, the operation count of Proposition~\ref{prop:C}, the reading of root multiplicities as ramification data, the point at infinity and the fibre over~$1$ in Theorem~\ref{thm:B}, the construction of the Galois rings and their cardinalities (so the count $(q-2)q^{k-1}$ in Theorem~\ref{thm:R}(1) and the value $2q^{k-1}=2\,\#\ker\pi$ in part~(2) rest on the text), and the bounds quoted from \cite{RS}. The correspondence between statements and declarations, the build and checker records, and a blueprint accompany the code.

\section*{Funding}
This research did not receive any specific grant from funding agencies in the public, commercial, or not-for-profit sectors.

\section*{Data availability}
The code, logs, expected outputs and manifests of the finite computations (the folder described above, with the entry point \texttt{run\_quick.sh}), together with the Lean sources, their build and checker records and the blueprint, are deposited at Zenodo, doi:10.5281/zenodo.23232181 (version 1.0.0); they also accompany the submission as supplementary material.

\section*{Declaration of generative AI and AI-assisted technologies in the manuscript preparation process}
During the preparation of this work the author used Claude and Claude Code (Anthropic) and ChatGPT (OpenAI) in order to draft, organise and edit the text and to prepare the literature survey; their use in the research is described in the section Methods and reproducibility. After using these tools, the author reviewed and edited the content as needed and takes full responsibility for the content of the published article.

\begin{thebibliography}{9}
\bibitem{BS} D.~Bartoli and P.~St\u{a}nic\u{a}, Reduced polynomial lifts of APN permutations over Galois rings and effective non-APN bounds, preprint, arXiv:2608.30808v1, 31 August 2026, doi:10.48550/arXiv.2608.30808.
\bibitem{BL} C.~Beierle and G.~Leander, New instances of quadratic APN functions, \emph{IEEE Trans.\ Inf.\ Theory} 68(1) (2022) 670--678, doi:10.1109/TIT.2021.3120698; arXiv:2009.07204v4.
\bibitem{Chen} R.~Chen, A general approach to permutation polynomials from quadratic forms, preprint, arXiv:2506.24012v1, 2025, doi:10.48550/arXiv.2506.24012.
\bibitem{TIT} H.~Fukui, Optimal reduced representatives of quadratic APN permutations in odd dimension, submitted to \emph{IEEE Trans.\ Inf.\ Theory} (IT-26-1499); preprint, Zenodo, doi:10.5281/zenodo.23074611.
\bibitem{KZ} C.~Kaspers and Y.~Zhou, A lower bound on the number of inequivalent APN functions, \emph{J.\ Combin.\ Theory Ser.~A} 186 (2022) 105554, doi:10.1016/j.jcta.2021.105554; cited by the numbering of arXiv:2002.00673v2.
\bibitem{LW} Y.~Li and M.~Wang, On EA-equivalence of certain permutations to power mappings, \emph{Des.\ Codes Cryptogr.} 58(3) (2011) 259--269, doi:10.1007/s10623-010-9406-8.
\bibitem{RS} S.~R{\o}njom and A.~Sandrib, On the differential uniformity of polynomials over Galois rings, \emph{Cryptogr.\ Commun.} 18(6) (2026) 2057--2077, doi:10.1007/s12095-026-00895-x.
\bibitem{Stone} D.~Stone, \texttt{apn-critical-point-counterexample}, GitHub repository, created 27 September 2026, \url{https://github.com/drewstone/apn-critical-point-counterexample} (accessed 7 October 2026).
\bibitem{YKBL} Y.~Yu, N.~Kaleyski, L.~Budaghyan and Y.~Li, Classification of quadratic APN functions with coefficients in $\F_2$ for dimensions up to 9, \emph{Finite Fields Appl.} 68 (2020) 101733, doi:10.1016/j.ffa.2020.101733.
\end{thebibliography}
\end{document}
